% =============================================================================
\chapter{Q-Learning et SARSA}
\label{ch:q_learning}
% =============================================================================

En 1989, Christopher Watkins, dans sa th\`ese de doctorat \`a Cambridge, propose un algorithme d'une simplicit\'e trompeuse~: \`a chaque transition, mettre \`a jour la valeur d'une paire \'etat-action en utilisant le maximum sur les actions suivantes. C'est le \emph{Q-learning}, le premier algorithme d'apprentissage par renforcement dont on a pu prouver la convergence vers la politique optimale \emph{sans} conna\^itre le mod\`ele de l'environnement. Presque simultan\'ement, Rummery et Niranjan proposent \emph{SARSA}, une variante \emph{on-policy}. La distinction entre ces deux approches --- \emph{off-policy} et \emph{on-policy} --- est l'une des lignes de fracture fondamentales de l'apprentissage par renforcement.

\begin{intuition}
Q-Learning et SARSA sont les algorithmes fondamentaux de contr\^ole sans mod\`ele.
Ils apprennent directement la fonction $Q(s,a)$ \`a partir de l'exp\'erience.
SARSA est \emph{on-policy} (en-politique)~: il apprend la valeur de la politique
qu'il suit. Q-Learning est \emph{off-policy} (hors-politique)~: il apprend la
valeur de la politique optimale ind\'ependamment de la politique d'exploration.
\end{intuition}

% -----------------------------------------------------------------------------
\section{Contr\^ole TD --- Cadre g\'en\'eral}
% -----------------------------------------------------------------------------

Le probl\`eme de \textbf{contr\^ole} consiste \`a trouver une politique optimale
$\pi^*$ sans conna\^itre le mod\`ele. On combine l'estimation de $Q^\pi$ par
m\'ethodes TD avec l'am\'elioration $\epsilon$-gloutonne, dans l'esprit de GPI.

% -----------------------------------------------------------------------------
\section{SARSA --- Contr\^ole en-politique}
% -----------------------------------------------------------------------------

\begin{definition}[SARSA]
L'algorithme SARSA met \`a jour $Q(S_t, A_t)$ en utilisant le quintuplet
$(S_t, A_t, R_{t+1}, S_{t+1}, A_{t+1})$~:
\[
    Q(S_t, A_t) \leftarrow Q(S_t, A_t) + \alpha \big[
    R_{t+1} + \gamma Q(S_{t+1}, A_{t+1}) - Q(S_t, A_t) \big].
\]
\end{definition}

\begin{formules}[title=\textbf{Mise \`a jour SARSA}]
\[
    Q(S_t, A_t) \leftarrow Q(S_t, A_t) + \alpha \big[
    R_{t+1} + \gamma Q(S_{t+1}, A_{t+1}) - Q(S_t, A_t) \big]
\]
o\`u $A_{t+1} \sim \pi_\epsilon(\cdot | S_{t+1})$ est l'action effectivement
choisie par la politique $\epsilon$-gloutonne.
\end{formules}

\begin{algorithme}[title=\textbf{SARSA}]
\begin{enumerate}
    \item Initialiser $Q(s,a)$ arbitrairement, $Q(\text{terminal}, \cdot) = 0$
    \item Pour chaque \'episode~:
    \begin{enumerate}
        \item Initialiser $S_0$~; choisir $A_0$ selon $\pi_\epsilon$ d\'eriv\'ee de $Q$
        \item Pour chaque pas $t = 0, 1, \ldots$ jusqu'\`a terminaison~:
        \begin{enumerate}
            \item Ex\'ecuter $A_t$, observer $R_{t+1}, S_{t+1}$
            \item Choisir $A_{t+1}$ selon $\pi_\epsilon$ d\'eriv\'ee de $Q$
            \item $Q(S_t, A_t) \leftarrow Q(S_t, A_t) + \alpha[R_{t+1} + \gamma Q(S_{t+1}, A_{t+1}) - Q(S_t, A_t)]$
        \end{enumerate}
    \end{enumerate}
\end{enumerate}
\end{algorithme}

\begin{theoreme}[Convergence de SARSA]
Sous les conditions de Robbins-Monro et si tous les couples $(s,a)$ sont visit\'es
infiniment souvent (condition GLIE~: Greedy in the Limit with Infinite Exploration),
SARSA converge vers $Q^*$ presque s\^urement.
\end{theoreme}

% -----------------------------------------------------------------------------
\section{Q-Learning --- Contr\^ole hors-politique}
% -----------------------------------------------------------------------------

\begin{definition}[Q-Learning]
L'algorithme Q-Learning de Watkins (1989) met \`a jour~:
\[
    Q(S_t, A_t) \leftarrow Q(S_t, A_t) + \alpha \big[
    R_{t+1} + \gamma \max_{a'} Q(S_{t+1}, a') - Q(S_t, A_t) \big].
\]
\end{definition}

\begin{formules}[title=\textbf{Mise \`a jour Q-Learning}]
\[
    Q(S_t, A_t) \leftarrow Q(S_t, A_t) + \alpha \Big[
    R_{t+1} + \gamma \max_{a' \in \mathcal{A}} Q(S_{t+1}, a') - Q(S_t, A_t) \Big]
\]
\end{formules}

La diff\'erence cl\'e avec SARSA est l'utilisation du $\max$ au lieu de $Q(S_{t+1}, A_{t+1})$.
Q-Learning apprend $Q^*$ directement, quelle que soit la politique de comportement
(tant que tous les couples $(s,a)$ sont visit\'es).

\begin{algorithme}[title=\textbf{Q-Learning}]
\begin{enumerate}
    \item Initialiser $Q(s,a)$ arbitrairement, $Q(\text{terminal}, \cdot) = 0$
    \item Pour chaque \'episode~:
    \begin{enumerate}
        \item Initialiser $S_0$
        \item Pour chaque pas $t = 0, 1, \ldots$ jusqu'\`a terminaison~:
        \begin{enumerate}
            \item Choisir $A_t$ selon $\pi_\epsilon$ d\'eriv\'ee de $Q$
            \item Ex\'ecuter $A_t$, observer $R_{t+1}, S_{t+1}$
            \item $Q(S_t, A_t) \leftarrow Q(S_t, A_t) + \alpha[R_{t+1} + \gamma \max_{a'} Q(S_{t+1}, a') - Q(S_t, A_t)]$
        \end{enumerate}
    \end{enumerate}
\end{enumerate}
\end{algorithme}

\begin{theoreme}[Convergence de Q-Learning]
\label{thm:q_learning_conv}
Si chaque couple $(s,a)$ est visit\'e infiniment souvent et si les taux
d'apprentissage $\alpha_t(s,a)$ satisfont les conditions de Robbins-Monro~:
\[
    \sum_t \alpha_t(s,a) = \infty, \qquad \sum_t \alpha_t^2(s,a) < \infty,
\]
alors $Q(s,a) \to Q^*(s,a)$ presque s\^urement pour tout $(s,a)$.
\end{theoreme}

\begin{proof}[Id\'ee de la preuve]
La preuve repose sur la th\'eorie des approximations stochastiques.
On \'ecrit la mise \`a jour sous la forme~:
\[
    Q_{t+1}(s,a) = (1-\alpha_t)Q_t(s,a) + \alpha_t [R_{t+1} + \gamma \max_{a'} Q_t(S_{t+1}, a')].
\]
C'est un processus d'approximation stochastique de l'op\'erateur de Bellman
$\mathcal{T}^*$. La contraction de $\mathcal{T}^*$ en norme infinie garantit
la convergence sous les conditions de Robbins-Monro (th\'eor\`eme de Jaakkola
et al., 1994).
\end{proof}

% -----------------------------------------------------------------------------
\section{SARSA vs Q-Learning~: l'exemple de la falaise}
% -----------------------------------------------------------------------------

\begin{exemple}[Cliff Walking]
L'environnement \texttt{CliffWalking-v0} illustre la diff\'erence~:
\begin{itemize}
    \item \textbf{Q-Learning} apprend la politique optimale (le long de la falaise),
    mais la politique d'exploration peut faire tomber l'agent.
    \item \textbf{SARSA} apprend une politique plus s\^ure (loin de la falaise)
    car il tient compte du comportement $\epsilon$-glouton.
\end{itemize}
Q-Learning optimise le retour de la politique \emph{cible} (gloutonne), tandis que
SARSA optimise le retour de la politique \emph{effective} ($\epsilon$-gloutonne).
\end{exemple}

% -----------------------------------------------------------------------------
\section{Expected SARSA}
% -----------------------------------------------------------------------------

\begin{definition}[Expected SARSA]
\textbf{Expected SARSA} utilise l'esp\'erance sur les actions futures au lieu
d'une action \'echantillonn\'ee~:
\[
    Q(S_t, A_t) \leftarrow Q(S_t, A_t) + \alpha \Big[
    R_{t+1} + \gamma \sum_{a'} \pi(a' | S_{t+1}) Q(S_{t+1}, a') - Q(S_t, A_t) \Big].
\]
\end{definition}

\begin{remarque}
Expected SARSA g\'en\'eralise \`a la fois SARSA et Q-Learning~:
\begin{itemize}
    \item Si $\pi$ est la politique $\epsilon$-gloutonne courante, c'est une
    version \`a variance r\'eduite de SARSA.
    \item Si $\pi$ est la politique gloutonne, c'est exactement Q-Learning.
\end{itemize}
\end{remarque}

% -----------------------------------------------------------------------------
\section{Double Q-Learning}
% -----------------------------------------------------------------------------

\begin{attention}[title=\textbf{Biais de maximisation}]
Q-Learning souffre d'un \textbf{biais de maximisation}~: $\E[\max_a Q(s,a)] \geq
\max_a \E[Q(s,a)]$ (in\'egalit\'e de Jensen). Ce biais conduit \`a surestimer
syst\'ematiquement les valeurs Q.
\end{attention}

\begin{definition}[Double Q-Learning]
\textbf{Double Q-Learning} (Hasselt, 2010) maintient deux tables $Q^A$ et $Q^B$.
\`A chaque pas, avec probabilit\'e $0.5$~:
\[
    Q^A(S_t, A_t) \leftarrow Q^A(S_t, A_t) + \alpha \Big[
    R_{t+1} + \gamma Q^B\big(S_{t+1}, \argmax_{a'} Q^A(S_{t+1}, a')\big) - Q^A(S_t, A_t) \Big]
\]
ou on met \`a jour $Q^B$ sym\'etriquement. L'action est s\'electionn\'ee par une
table et \'evalu\'ee par l'autre, \'eliminant le biais.
\end{definition}

% -----------------------------------------------------------------------------
\section{SARSA($\lambda$) et Q($\lambda$)}
% -----------------------------------------------------------------------------

\begin{definition}[SARSA($\lambda$)]
SARSA($\lambda$) combine SARSA avec les traces d'\'eligibilit\'e~:
\[
    e_t(s,a) = \gamma \lambda \, e_{t-1}(s,a) + \mathbb{1}_{S_t=s, A_t=a}
\]
\[
    Q(s,a) \leftarrow Q(s,a) + \alpha \, \delta_t \, e_t(s,a)
\]
o\`u $\delta_t = R_{t+1} + \gamma Q(S_{t+1}, A_{t+1}) - Q(S_t, A_t)$.
\end{definition}

% -----------------------------------------------------------------------------
\section{Impl\'ementation en Python}
% -----------------------------------------------------------------------------

\begin{codeblock}[title=\textbf{Q-Learning tabulaire}]
\begin{minted}{python}
import numpy as np
import gymnasium as gym

def q_learning(env, n_episodes=5000, alpha=0.1, gamma=0.99,
               epsilon=0.1):
    """Q-Learning tabulaire avec politique epsilon-gloutonne."""
    nS = env.observation_space.n
    nA = env.action_space.n
    Q = np.zeros((nS, nA))

    def epsilon_greedy(s):
        if np.random.random() < epsilon:
            return env.action_space.sample()
        return np.argmax(Q[s])

    rewards_per_episode = []
    for ep in range(n_episodes):
        s, _ = env.reset()
        total_reward = 0
        done = False
        while not done:
            a = epsilon_greedy(s)
            s_next, r, terminated, truncated, _ = env.step(a)
            done = terminated or truncated
            # Mise a jour Q-Learning
            target = r + gamma * np.max(Q[s_next]) * (not terminated)
            Q[s, a] += alpha * (target - Q[s, a])
            s = s_next
            total_reward += r
        rewards_per_episode.append(total_reward)

    pi = np.argmax(Q, axis=1)
    return Q, pi, rewards_per_episode
\end{minted}
\end{codeblock}

\begin{codeblock}[title=\textbf{SARSA tabulaire}]
\begin{minted}{python}
def sarsa(env, n_episodes=5000, alpha=0.1, gamma=0.99,
          epsilon=0.1):
    """SARSA tabulaire avec politique epsilon-gloutonne."""
    nS = env.observation_space.n
    nA = env.action_space.n
    Q = np.zeros((nS, nA))

    def epsilon_greedy(s):
        if np.random.random() < epsilon:
            return env.action_space.sample()
        return np.argmax(Q[s])

    rewards_per_episode = []
    for ep in range(n_episodes):
        s, _ = env.reset()
        a = epsilon_greedy(s)
        total_reward = 0
        done = False
        while not done:
            s_next, r, terminated, truncated, _ = env.step(a)
            done = terminated or truncated
            a_next = epsilon_greedy(s_next) if not done else 0
            # Mise a jour SARSA
            target = r + gamma * Q[s_next, a_next] * (not terminated)
            Q[s, a] += alpha * (target - Q[s, a])
            s, a = s_next, a_next
            total_reward += r
        rewards_per_episode.append(total_reward)

    pi = np.argmax(Q, axis=1)
    return Q, pi, rewards_per_episode
\end{minted}
\end{codeblock}

\begin{codeblock}[title=\textbf{Comparaison sur CliffWalking}]
\begin{minted}{python}
import matplotlib.pyplot as plt

env = gym.make("CliffWalking-v0")

Q_ql, pi_ql, rewards_ql = q_learning(env, n_episodes=5000)
Q_sa, pi_sa, rewards_sa = sarsa(env, n_episodes=5000)

# Lissage par moyenne glissante
def smooth(x, window=100):
    return np.convolve(x, np.ones(window)/window, mode='valid')

plt.figure(figsize=(10, 5))
plt.plot(smooth(rewards_ql), label='Q-Learning')
plt.plot(smooth(rewards_sa), label='SARSA')
plt.xlabel('Episode')
plt.ylabel('Recompense cumulee (lissee)')
plt.legend()
plt.title('CliffWalking: Q-Learning vs SARSA')
plt.show()
\end{minted}
\end{codeblock}

\begin{codeblock}[title=\textbf{Double Q-Learning}]
\begin{minted}{python}
def double_q_learning(env, n_episodes=5000, alpha=0.1,
                      gamma=0.99, epsilon=0.1):
    """Double Q-Learning tabulaire."""
    nS = env.observation_space.n
    nA = env.action_space.n
    QA = np.zeros((nS, nA))
    QB = np.zeros((nS, nA))

    def epsilon_greedy(s):
        if np.random.random() < epsilon:
            return env.action_space.sample()
        return np.argmax(QA[s] + QB[s])

    for ep in range(n_episodes):
        s, _ = env.reset()
        done = False
        while not done:
            a = epsilon_greedy(s)
            s_next, r, terminated, truncated, _ = env.step(a)
            done = terminated or truncated
            if np.random.random() < 0.5:
                a_star = np.argmax(QA[s_next])
                target = r + gamma * QB[s_next, a_star] * (not terminated)
                QA[s, a] += alpha * (target - QA[s, a])
            else:
                a_star = np.argmax(QB[s_next])
                target = r + gamma * QA[s_next, a_star] * (not terminated)
                QB[s, a] += alpha * (target - QB[s, a])
            s = s_next

    Q = (QA + QB) / 2
    pi = np.argmax(Q, axis=1)
    return Q, pi
\end{minted}
\end{codeblock}

% -----------------------------------------------------------------------------
\section{Exercices}
% -----------------------------------------------------------------------------

\begin{exercice}[Biais de maximisation]
Construisez un MDP \`a 2 \'etats o\`u Q-Learning surestime syst\'ematiquement
les valeurs. Ex\'ecutez 1000 r\'ep\'etitions et comparez les estimations de
Q-Learning, Double Q-Learning et la vraie valeur $Q^*$.
\end{exercice}

\begin{exercice}[Taux d'apprentissage]
\'Etudiez l'effet de $\alpha$ sur la convergence de Q-Learning pour
$\alpha \in \{0.01, 0.05, 0.1, 0.5, 1.0\}$ sur \texttt{Taxi-v3}.
Tracez les courbes de convergence.
\end{exercice}

\begin{exercice}[Expected SARSA]
Impl\'ementez Expected SARSA et montrez qu'il r\'eduit la variance par rapport
\`a SARSA sur \texttt{CliffWalking-v0}. Comparez les \'ecarts-types des
r\'ecompenses cumulatives sur 100 ex\'ecutions.
\end{exercice}

\begin{exercice}[Preuve formelle]
Montrez que Q-Learning est un cas sp\'ecial d'Expected SARSA lorsque la
politique cible est la politique gloutonne. \'Ecrivez la preuve formelle.
\end{exercice}
